\documentclass[11pt]{article}
\usepackage[margin=1in]{geometry}
\usepackage{amsmath,amssymb,amsthm}
\usepackage{graphicx}
\usepackage[hidelinks]{hyperref}
\emergencystretch=2em

\newtheorem{theorem}{Theorem}
\newtheorem{lemma}{Lemma}
\newtheorem{remark}{Remark}

\newcommand{\N}{\mathbb N}
\newcommand{\R}{\mathbb R}
\newcommand{\cnull}{c_0}
\newcommand{\Range}{\mathcal R}

\title{Partial Result: Two \(c_0\) Vector-Measure Range Subcases}
\author{Automatic Math Discovery Run \texttt{fa\_banach\_001}}
\date{June 15, 2026}

\begin{document}
\maketitle

\section*{Source Problem}

The source is Jose Rodriguez, \emph{Open problems in Banach spaces and
measure theory}. The arXiv identifier is 1607.07787.

Problem 4.2 appears on PDF page 12 and is attributed there to Anantharaman
and Diestel:
\begin{quote}
Let \(K\) be a weakly compact subset of \(c_0\). Is \(K\) contained in the
range of a \(c_0\)-valued vector measure?
\end{quote}

\begin{center}
\includegraphics[width=\textwidth]{figures/open_problem_crop.png}
\end{center}

\section*{Definitions}

For a countably additive vector measure \(\nu:\Sigma\to X\), its range is
\[
\Range(\nu)=\{\nu(A): A\in\Sigma\}.
\]
Here countable additivity is meant in the norm of the target Banach space.
The problem asks whether every weakly compact \(K\subseteq c_0\) is contained
in \(\Range(\nu)\) for some \(c_0\)-valued vector measure \(\nu\).

For \(a=(a_n)\in c_0\) with \(a_n\ge 0\), write
\[
Q(a)=\{x\in c_0: |x_n|\le a_n\text{ for every }n\}.
\]

\section*{Idea of the Proof}

The coordinate-box case is realized one coordinate at a time. Put two
atomless probability intervals behind each coordinate and let the \(n\)-th
coordinate of the vector measure record a signed amount of mass from those
two intervals, scaled by \(a_n\). The tail condition \(a_n\to0\) supplies the
uniform \(c_0\)-control needed for countable additivity.

The second construction shows that the first theorem is not merely a compact
set observation. The set \(\{0,e_1,e_2,\ldots\}\) is weakly compact in
\(c_0\) but not norm compact. Rademacher functions give a vector measure
whose value on the half-space \(\{r_n=1\}\) is exactly \(e_n\).

\begin{theorem}[coordinate boxes]
Let \(a=(a_n)\in c_0\) with \(a_n\ge0\). Then \(Q(a)\) is contained in the
range of a countably additive \(c_0\)-valued vector measure. Consequently,
every norm-compact subset of \(c_0\) is contained in such a range.
\end{theorem}

\begin{proof}
Let
\[
\Omega=\bigsqcup_{n\ge1}(I_n^+\sqcup I_n^-),
\]
where each \(I_n^+\) and \(I_n^-\) is a copy of \([0,1]\) with Lebesgue
measure. Let \(\Sigma\) be the sigma algebra obtained from the disjoint union.
For \(A\in\Sigma\), define
\[
\nu(A)_n=a_n\bigl(\lambda(A\cap I_n^+)-\lambda(A\cap I_n^-)\bigr).
\]
Then \(|\nu(A)_n|\le a_n\), so \(\nu(A)\in c_0\).

We check countable additivity in \(c_0\). Let \(A_1,A_2,\ldots\) be pairwise
disjoint and put \(B_m=\bigcup_{j>m}A_j\). Since each coordinate is a finite
scalar signed measure, for every fixed \(n\) we have \(\nu(B_m)_n\to0\). Given
\(\varepsilon>0\), choose \(N\) such that \(\sup_{n>N}a_n<\varepsilon\). For
the finitely many coordinates \(n\le N\), choose \(m\) so large that
\(|\nu(B_m)_n|<\varepsilon\) for all \(n\le N\). Then
\(\|\nu(B_m)\|_\infty\le\varepsilon\). Hence the finite partial sums converge
in \(c_0\) to \(\nu(\bigcup_jA_j)\).

Now fix \(x\in Q(a)\). For each \(n\) with \(a_n>0\), choose a measurable
set \(B_n\subseteq I_n^+\) of measure \(x_n/a_n\) if \(x_n\ge0\), and choose
\(B_n\subseteq I_n^-\) of measure \(-x_n/a_n\) if \(x_n<0\). If \(a_n=0\),
then necessarily \(x_n=0\), and we take \(B_n=\varnothing\). With
\(B=\bigcup_nB_n\), the definition gives \(\nu(B)=x\). Thus
\(Q(a)\subseteq\Range(\nu)\).

Finally, if \(K\subseteq c_0\) is norm compact, set
\[
a_n=\sup_{x\in K}|x_n|.
\]
Compactness of \(K\) implies \(a_n\to0\): otherwise a finite epsilon-net in
\(c_0\) would give uniformly small tails for all points of \(K\), a
contradiction. Therefore \(K\subseteq Q(a)\), and the first part applies.
\end{proof}

\begin{theorem}[the canonical weakly compact sequence]
The set \(\{0\}\cup\{e_n:n\in\N\}\subseteq c_0\) is contained in the range of
a countably additive \(c_0\)-valued vector measure.
\end{theorem}

\begin{proof}
Let \((r_n)_{n\ge1}\) be the standard Rademacher functions on \([0,1]\). For
Lebesgue measurable \(A\subseteq[0,1]\), define
\[
\nu(A)=\bigl(2\int_A r_n(t)\,dt\bigr)_{n\ge1}.
\]
We first show that \(\nu(A)\in c_0\). Approximate \(1_A\) in \(L^1[0,1]\) by
a dyadic step function \(s\) with mesh \(2^{-m}\). For every \(n>m\),
\(\int_0^1 s(t)r_n(t)\,dt=0\). Hence
\[
\left|\int_A r_n(t)\,dt\right|
=\left|\int_0^1(1_A(t)-s(t))r_n(t)\,dt\right|
\le \|1_A-s\|_1
\]
for all \(n>m\). Since the approximation error can be arbitrarily small,
\(\int_A r_n\to0\), so \(\nu(A)\in c_0\).

Countable additivity is also in the norm of \(c_0\). If
\((A_j)_{j\ge1}\) are pairwise disjoint and \(C_m=\bigcup_{j>m}A_j\), then
\(\lambda(C_m)\to0\). Thus
\[
\|\nu(C_m)\|_\infty
\le 2\lambda(C_m)\to0,
\]
which is exactly the norm convergence of the tails of the countable sum.

Let \(E_n=\{t:r_n(t)=1\}\). Orthogonality and independence of the Rademacher
functions give
\[
2\int_{E_n}r_j(t)\,dt=\delta_{jn}.
\]
Therefore \(\nu(E_n)=e_n\), while \(\nu(\varnothing)=0\). This proves that
\(\{0\}\cup\{e_n:n\in\N\}\subseteq\Range(\nu)\).
\end{proof}

\begin{remark}
The set \(\{0\}\cup\{e_n:n\in\N\}\) is weakly compact in \(c_0\), since
\(e_n\to0\) weakly against the dual space \(\ell_1\). It is not norm compact,
because \(\|e_n-e_m\|_\infty=1\) for \(n\ne m\). Thus the Rademacher theorem
is a genuine non-norm-compact positive instance of the source problem.
\end{remark}

\section*{Scope and Limitations}

This packet is a partial result only. It proves the Anantharaman-Diestel
question for norm-compact subsets of \(c_0\), and it gives one basic weakly
compact non-norm-compact example. It does not settle arbitrary weakly compact
subsets of \(c_0\), where no coordinatewise null majorant need exist and where
the weak compactness may encode substantially more complicated geometry than
the single weakly null basis sequence.

\section*{Novelty and Search Notes}

The run indexes were searched for arXiv:1607.07787, for the phrases
``Anantharaman-Diestel'', ``weakly compact subset of c0'', ``vector measure'',
``range of a c0-valued vector measure'', and for nearby Banach-space measure
theory keywords. No existing run packet for Problem 4.2 or for these two
subcases was found. The source paper itself presents Problem 4.2 as open.

This result is constructive and elementary. No later-paper identification is
being claimed; the packet should be reviewed as a newly generated partial
answer to the extracted problem.

\section*{Human Review Recommendation}

Review the two norm-convergence arguments for countable additivity. In the
coordinate-box construction, the only uniformity input is \(a_n\to0\). In the
Rademacher construction, the only uniformity input for additivity is
\(\lambda(C_m)\to0\) for residual unions in a finite measure space. The
containment computations are coordinatewise.

\begin{thebibliography}{9}
\bibitem{Rodriguez}
J. Rodriguez,
\emph{Open problems in Banach spaces and measure theory},
arXiv:1607.07787.

\bibitem{AnantharamanDiestel}
R. Anantharaman and J. Diestel,
The problem cited as the source of Problem 4.2 in Rodriguez,
\emph{Open problems in Banach spaces and measure theory}.
\end{thebibliography}

\end{document}
